\documentclass[11pt]{article}
\usepackage[T1]{fontenc}
\usepackage{textcomp}
\usepackage[margin=0.75in]{geometry}
\usepackage{amsmath,amssymb,amsthm}
\usepackage{array,booktabs}
\usepackage{hyperref}
\setlength{\parindent}{0pt}
\newtheorem{theorem}{Theorem}
\theoremstyle{definition}
\newtheorem{definition}{Definition}
\title{Flow Proof Artifact: Ratio-componendo-dividendo}
\author{Generated by \texttt{flow doc proof}}
\date{Flow Proof Book}
\begin{document}
\maketitle
Proportional magnitudes satisfy componendo and dividendo.\\[0.5em]
\textbf{Source.} euclid — Elements, Book V, Proposition 13\\[0.5em]
\bigskip
\noindent{\large \textbf{Derived fact 1.} \textit{a:b = c:d implies (a+b):(a-b) = (c+d):(c-d)} \hfill Ratio \cdot proportion \cdot componendo and dividendo hold}\par
\smallskip\noindent\textit{Coordinate: Ratio \cdot proportion \cdot componendo and dividendo hold}\par
\smallskip\noindent\textit{Source: euclid}\par
\smallskip\noindent\textit{Built on: componendo holds, for proportion on Ratio, dividendo holds, for proportion on Ratio}\par
\medskip
\noindent\textbf{Goal.} a:b = c:d implies (a+b):(a-b) = (c+d):(c-d)\par
\noindent$\displaystyle \forall a \in Magnitude \forall b \in Magnitude \forall c \in Magnitude \forall d \in Magnitude\quad ratio(a + b, a - b) = ratio(c + d, c - d)$\par
\medskip
\noindent
\begin{tabular}{@{} >{\raggedright\arraybackslash}p{0.06\textwidth} >{\raggedright\arraybackslash}p{0.47\textwidth} >{\raggedleft\arraybackslash}p{0.41\textwidth} @{}}
\textbf{\#} & \textbf{Proof} & \textbf{Mathematics} \\
\midrule
\textbf{1.} & We prove that componendo and dividendo hold for proportion on Ratio. & \\[0.45em]
\textbf{2.} & We invoke the derived fact governing proportion on Ratio: componendo holds, for proportion on Ratio (instantiated for a, b, c, d). & \\[0.45em]
\textbf{3.} & We invoke the derived fact governing proportion on Ratio: dividendo holds, for proportion on Ratio (instantiated for a, b, c, d). & \\[0.45em]
\textbf{4.} & From step 2 and step 3, this implies ratio(a plus b, a - b) equals ratio(c plus d, c - d). Hence proven. & $\displaystyle ratio(a + b, a - b) = ratio(c + d, c - d)$ \\[0.45em]
\bottomrule
\end{tabular}
\label{thm:Ratio-proportion-componendo_and_dividendo_hold}
\par\medskip\hrule
\end{document}
